\documentclass[
%toc=listof, % Abbildungsverzeichnis sowie Tabellenverzeichnis in das Inhaltsverzeichnis aufnehmen
    toc=bibliography, % Literaturverzeichnis in das Inhaltsverzeichnis aufnehmen
    footnotes=multiple, % Trennen von direkt aufeinander folgenden Fußnoten
    parskip=half % vertikalen Abstand zwischen Absätzen verwenden anstatt horizontale Einrückung von Folgeabsätzen
]{scrartcl}

% ----- settings --------------------------------
\usepackage[utf8]{inputenc} % erlaubt Sonderzeichen im source; hier, weil meta es bereits braucht und packages.tex meta.tex nutzt
\newcommand{\titel}{Erweiterung einer Webanwendung um Task Notifications}

\newcommand{\autor}{Björn Nowak}
\newcommand{\kompletterTitel}{\titel}

\newcommand{\pruefungstermin}{Sommer 2015}
\newcommand{\ort}{Oldenburg}
\newcommand{\abgabetermin}{06.05.2015}
\newcommand{\ausbildungsberuf}{Fachinformatiker Anwendungsentwicklung}
\newcommand{\betreff}{Dokumentation zur betrieblichen Projektarbeit}

\newcommand{\copyrightCompany}{2015 BTC AG}

% Anpassung an Landessprache ---------------------------------------------------
\usepackage[ngerman]{babel}

% Umlaute ----------------------------------------------------------------------
%   Umlaute/Sonderzeichen wie äüöß direkt im Quelltext verwenden (CodePage).
% ------------------------------------------------------------------------------
\usepackage[T1]{fontenc} % Erlaubt automatische Trennung von Worten mit Umlauten.
\usepackage{textcomp} % Euro-Zeichen etc.

% Schrift ----------------------------------------------------------------------
\usepackage{lmodern} % bessere Fonts
\usepackage{relsize} % Schriftgröße relativ festlegen

% Tabellen ---------------------------------------------------------------------
\PassOptionsToPackage{table}{xcolor}
\usepackage{tabularx}
\usepackage{multirow}
\usepackage{bigstrut}
% für lange Tabellen
\usepackage{longtable}
\usepackage{array}
\usepackage{ragged2e}
\usepackage{lscape}
\newcolumntype{w}[1]{>{\raggedleft\hspace{0pt}}p{#1}} % Spaltendefinition rechtsbündig mit definierter Breite

% Grafiken ---------------------------------------------------------------------
\usepackage[dvips,final]{graphicx} % Einbinden von JPG-Grafiken ermöglichen
\usepackage{graphics} % keepaspectratio
\usepackage{floatflt} % zum Umfließen von Bildern
\graphicspath{{Bilder/}} % hier liegen die Bilder des Dokuments

\usepackage{tikz} % visualisation
\usetikzlibrary{trees} % directory trees

\usepackage{wrapfig}

% Sonstiges --------------------------------------------------------------------
\usepackage{amssymb} % fuer /checkmark
\usepackage{amsmath,amsfonts} % Befehle aus AMSTeX für mathematische Symbole
\usepackage{enumitem} % anpassbare Enumerates/Itemizes
\usepackage{xspace} % sorgt dafür, dass Leerzeichen hinter parameterlosen Makros nicht als Makroendezeichen interpretiert werden
\usepackage{xparse} % fuer \NewDocumentCommand, ermoeglicht mehrere optionale args
\usepackage{makeidx} % für Index-Ausgabe mit \printindex
\usepackage{acronym} % es werden alle Definitionen aufgelistet | nur benutze : \usepackage[printonlyused]{acronym}

% Einfache Definition der Zeilenabstände und Seitenränder etc.
\usepackage{setspace}
\usepackage{geometry}
% TODO may include
%\usepackage[compact]{titlesec} % Zeilenabstände

% Symbolverzeichnis
\usepackage[intoc]{nomencl}
\let\abbrev\nomenclature
\renewcommand{\nomname}{Abkürzungsverzeichnis}
\setlength{\nomlabelwidth}{.25\hsize}
\renewcommand{\nomlabel}[1]{#1 \dotfill}
\setlength{\nomitemsep}{-\parsep}

\usepackage{url} % URL verlinken, lange URLs umbrechen etc.
\usepackage{chngcntr} % fortlaufendes Durchnummerieren der Fußnoten
\usepackage{ifthen} % bei der Definition eigener Befehle benötigt
\usepackage{todonotes} % definiert u.a. die Befehle \todo und \listoftodos

\usepackage[
    backend=biber, %is default
    style=ieee, % numeric, alphabetic, authoryear, authortitle, science, ieee
    sorting=nyt
]{biblatex} % literaturverzeichnis
\addbibresource{template.bib}

% nty   Sort by name, title, year.
% nyt   Sort by name, year, title.
% nyvt  Sort by name, year, volume, title.
% anyt  Sort by alphabetic label, name, year, title.
% anyvt Sort by alphabetic label, name, year, volume, title.
% ynt   Sort by year, name, title.
% ydnt  Sort by year (descending), name, title.
% none  Do not sort at all. All entries are processed in citation order.


% PDF-Optionen -----------------------------------------------------------------
\usepackage{pdfpages}
\pdfminorversion=5 % erlaubt das Einfügen von pdf-Dateien bis Version 1.7, ohne eine Fehlermeldung zu werfen (keine Garantie für fehlerfreies Einbetten!)
\usepackage[
    bookmarks,
    bookmarksnumbered,
    bookmarksopen=true,
    bookmarksopenlevel=1,
    colorlinks=true,
% diese Farbdefinitionen zeichnen Links im PDF farblich aus
    linkcolor=MyBlau, % einfache interne Verknüpfungen
    anchorcolor=MyBlau,% Ankertext
    citecolor=MyBlau, % Verweise auf Literaturverzeichniseinträge im Text
    filecolor=MyBlau, % Verknüpfungen, die lokale Dateien öffnen
    menucolor=MyBlau, % Acrobat-Menüpunkte
    urlcolor=MyBlau,
% diese Farbdefinitionen sollten für den Druck verwendet werden (alles schwarz)
%linkcolor=black, % einfache interne Verknüpfungen
%anchorcolor=black, % Ankertext
%citecolor=black, % Verweise auf Literaturverzeichniseinträge im Text
%filecolor=black, % Verknüpfungen, die lokale Dateien öffnen
%menucolor=black, % Acrobat-Menüpunkte
%urlcolor=black,
%
%backref, % Quellen werden zurück auf ihre Zitate verlinkt
    pdftex,
    plainpages=false, % zur korrekten Erstellung der Bookmarks
    pdfpagelabels=true, % zur korrekten Erstellung der Bookmarks
    hypertexnames=false, % zur korrekten Erstellung der Bookmarks
    linktoc=all%linktocpage % Seitenzahlen anstatt Text im Inhaltsverzeichnis verlinken
]{hyperref}
% Befehle, die Umlaute ausgeben, führen zu Fehlern, wenn sie hyperref als Optionen übergeben werden
\hypersetup{
    pdftitle={\titel},
    pdfauthor={\autor},
    pdfcreator={\autor},
    pdfsubject={\titel},
    pdfkeywords={\titel},
}


% zum Einbinden von Programmcode -----------------------------------------------
\usepackage{listings}
\usepackage{xcolor}

\definecolor{DEFhellgelb}{rgb}{1,1,0.9}
\definecolor{DEFKeys}{rgb}{0,0,1}
\definecolor{DEFIdentifier}{rgb}{0,0,0}
\definecolor{DEFComments}{rgb}{0,0.5,0}
\definecolor{DEFString}{rgb}{1,0,0}

\definecolor{JAVAred}{rgb}{0.6,0,0}
\definecolor{JAVAgreen}{rgb}{0.25,0.5,0.35}
\definecolor{JAVApurple}{rgb}{0.5,0,0.35}
\definecolor{JAVAdocblue}{rgb}{0.25,0.35,0.75}

\definecolor{forestgreen}{RGB}{34,139,34}
\definecolor{orangered}{RGB}{239,134,64}
\definecolor{darkblue}{rgb}{0.0,0.0,0.6}
\definecolor{gray}{rgb}{0.4,0.4,0.4}

\definecolor{HTMLdarkgray}{rgb}{0.4, 0.4, 0.4}
\definecolor{HTMLocherCode}{rgb}{1, 0.5, 0} % #FF7F00 -> rgb(239, 169, 0)
\definecolor{HTMLblueCode}{rgb}{0, 0, 0.93} % #0000EE -> rgb(0, 0, 238)
\definecolor{HTMLgreenCode}{rgb}{0, 0.6, 0} % #009900 -> rgb(0, 153, 0)

\lstset{
    float=hbp,
    basicstyle=\footnotesize,
    identifierstyle=\color{DEFIdentifier},
    keywordstyle=\color{DEFKeys},
    stringstyle=\color{DEFString},
    commentstyle=\color{DEFComments},
    backgroundcolor=\color{white},
    columns=flexible,
    tabsize=2,
    frame=single,
    extendedchars=true,
    showspaces=false,
    showstringspaces=false,
    numbers=left,
    numberstyle=\tiny,
    breaklines=true,
    breakautoindent=true,
    captionpos=b,
}

\lstdefinestyle{Java}{
    language=Java,
    tabsize=4,
    tabsize=4,
    breakatwhitespace=false,
    resetmargins=true
%
    basicstyle=\footnotesize\ttfamily,
    keywordstyle=\color{JAVApurple}\bfseries,
    stringstyle=\color{JAVAred},
    commentstyle=\color{JAVAgreen},
%
    sensitive=true,
    morecomment=[l]{//},
    morecomment=[s][\color{JAVAdocblue}]{/**}{*/},
    morestring=[b]",
    morestring=[b]',
    morekeywords={
        public,void,package,import,throws,
        abstract,and,array,as,break,case,catch,cfunction,class,clone,const,
        continue,declare,default,do,else,elseif,enddeclare,endfor,endforeach,
        endif,endswitch,endwhile,extends,final,for,foreach,function,global,
        goto,if,implements,interface,instanceof,namespace,new,old_function,or,
        private,protected,public,static,switch,throw,try,use,var,while,xor
        die,echo,empty,exit,eval,include,include_once,isset,list,require,
        require_once,return,print,unset},
}

\lstdefinestyle{XML} {
    basicstyle=\footnotesize,
    columns=flexible,
    tabsize=2,
    frame=single,
    extendedchars=true,
    showspaces=false,
    showstringspaces=false,
    numbers=left,
    numberstyle=\tiny,
    breaklines=true,
    breakautoindent=true,
    captionpos=b,
%
    language=XML,
    extendedchars=true,
    breaklines=true,
    breakatwhitespace=true,
    emph={},
    emphstyle=\color{red},
    basicstyle=\ttfamily,
    columns=fullflexible,
    commentstyle=\color{gray}\upshape,
    morestring=[b]",
    morecomment=[s]{<?}{?>},
    morecomment=[s][\color{forestgreen}]{<!--}{-->},
    keywordstyle=\color{orangered},
    stringstyle=\ttfamily\color{black}\normalfont,
    tagstyle=\color{darkblue}\textbf,
    morekeywords={attribute,xmlns,version,type,release},
}

\lstdefinelanguage{CSS} {
    morekeywords={
        color,background,margin,padding,font,weight,display,position,
        top,left,right,bottom,list,style,border,size,white,space,min,width},
    sensitive=false,
    morecomment=[l]{//},
    morecomment=[s]{/*}{*/},
    morestring=[b]',
    morestring=[b]"
}

\lstdefinelanguage{JavaScript} {
    showstringspaces=false,
    keywordstyle=\ttfamily\bfseries\color{darkblue},
    identifierstyle=\ttfamily,
    stringstyle=\ttfamily\color{forestgreen},
    commentstyle=\color{gray},
    rulecolor=\color{gray},
    xleftmargin=5pt,
    xrightmargin=5pt,
    aboveskip=\bigskipamount,
    belowskip=\bigskipamount,
%
    morekeywords={
        break,const,continue,delete,do,while,export,for,in,function,
        if,else,import,in,instanceOf,label,let,new,return,switch,this,
        throw,try,catch,typeof,var,void,with,yield
    },
    sensitive=false,
    morecomment=[l]{//},
    morecomment=[s]{/*}{*/},
    morestring=[b]",
    morestring=[d]'
}
% Seitenränder -----------------------------------------------------------------
\setlength{\topskip}{\ht\strutbox} % behebt Warnung von geometry
\geometry{a4paper,left=20mm,right=20mm,top=25mm,bottom=35mm}

\usepackage[
    automark, % Kapitelangaben in Kopfzeile automatisch erstellen
    headsepline, % Trennlinie unter Kopfzeile
    footsepline, % Trennlinie über Fußzeile
    ilines % Trennlinie linksbündig ausrichten
]{scrpage2}
% TODO old {scrpage2}
% TODO new {scrlayer-scrpage}

% Kopf- und Fußzeilen ----------------------------------------------------------
\pagestyle{scrheadings}
%\renewcommand{\chapterpagestyle}{scrheadings}
\clearscrheadfoot

% Kopfzeile
\renewcommand{\headfont}{\normalfont} % Schriftform der Kopfzeile
\ihead{\large{\textsc{\titel}} \\[2ex] \textit{\headmark}}
\chead{}
\ohead{\includegraphics[scale=0.125]{picture/btc_logo_rgb_office}}
\setlength{\headheight}{15mm} % Höhe der Kopfzeile
%\setheadwidth[0pt]{textwithmarginpar} % Kopfzeile über den Text hinaus verbreitern (falls Logo den Text überdeckt) 

% Fußzeile
\ifoot{\copyright\ \copyrightCompany\ - \autor}
\cfoot{}
\ofoot{\pagemark}
\setlength{\footskip}{1.2cm}

% Zeilenabstände
\singlespacing % Zeilenabstand 1,5 Zeilen
\frenchspacing % erzeugt ein wenig mehr Platz hinter einem Punkt 


% Schusterjungen und Hurenkinder vermeiden
\clubpenalty = 10000
\widowpenalty = 10000
\displaywidowpenalty = 10000

% Quellcode-Ausgabe formatieren
\lstset{numbers=left, numberstyle=\tiny, numbersep=5pt, breaklines=true}
\lstset{emph={square}, emphstyle=\color{red}, emph={[2]root,base}, emphstyle={[2]\color{blue}}}

\counterwithout{footnote}{section} % Fußnoten fortlaufend durchnummerieren
\setcounter{tocdepth}{2} % im Inhaltsverzeichnis werden die Kapitel bis zum Level der paragraphen übernommen
\setcounter{secnumdepth}{4} % Kapitel bis zum Level der paragraphen werden nummeriert
% \subsubsectionlevel = level 3 

% Aufzählungen anpassen
\renewcommand{\labelenumi}{\arabic{enumi}.}
\renewcommand{\labelenumii}{\arabic{enumi}.\arabic{enumii}.}
\renewcommand{\labelenumiii}{\arabic{enumi}.\arabic{enumii}.\arabic{enumiii}}

\renewcommand{\autodot}{} % Remove all end-of-counter dots

% Färbung:
\definecolor{MyBlau}{rgb}{0, 0.28, 0.56}

\definecolor{head}{RGB}{173,216,230} % {0.4, 0.6, 0.8}
\definecolor{sub}{RGB}{209,238,238}
\definecolor{odd}{RGB}{255,255,234}
\definecolor{hig}{RGB}{255,255,224}

% Abkürzungen, ggfs. mit korrektem Leerraum 
\newcommand{\bs}{$\backslash$\xspace}
\newcommand{\bspw}{bspw.\xspace}
\newcommand{\bzw}{bzw.\xspace}
\newcommand{\ca}{ca.\xspace}
\newcommand{\dahe}{\mbox{d.\,h.}\xspace}
\newcommand{\etc}{etc.\xspace}
\newcommand{\eur}[1]{\mbox{#1\,\texteuro}\xspace}
\newcommand{\evtl}{evtl.\xspace}
\newcommand{\ggfs}{ggf.\xspace}
\newcommand{\Ggfs}{Ggf.\xspace}
\newcommand{\gqq}[1]{\glqq{}#1\grqq{}}
\newcommand{\inkl}{inkl.\xspace}
\newcommand{\insb}{insb.\xspace}
\newcommand{\ua}{\mbox{u.\,a.}\xspace}
\newcommand{\usw}{usw.\xspace}
\newcommand{\Vgl}{Vgl.\xspace}
\newcommand{\zB}{\mbox{z.\,B.}\xspace}
\newcommand{\og}{\mbox{o.\,g.}\xspace}
\newcommand{\sog}{sog.\xspace}

% Befehle für häufig anfallende Aufgaben
\newcommand{\Abbildung}[1]{Abbildung~\ref{#1}}
\newcommand{\AbbildungV}[1]{Abbildung~\ref{#1}: \nameref{#1} auf S.~\pageref{#1}}
\newcommand{\Anhang}[1]{Anhang~\ref{#1}: \nameref{#1} auf S.~\pageref{#1}}
\newcommand{\AnhangK}[1]{Anhang~\ref{#1} auf S.~\pageref{#1}}
\newcommand{\AnhangKapitel}[1]{Kapitel~\ref{#1}: \nameref{#1} auf S.~\pageref{#1}}
\newcommand{\Tabelle}[1]{Tabelle~\ref{#1}: \nameref{#1}}
\newcommand{\TabelleK}[1]{Tabelle~\ref{#1} auf S.~\pageref{#1}}
\newcommand{\TabelleV}[1]{Tabelle~\ref{#1}: \nameref{#1} auf S.~\pageref{#1}}

\newcommand{\includegraphicsKeepAspectRatio}[2]{\includegraphics[width=#2\textwidth,height=#2\textheight,keepaspectratio]{#1}}
\newcommand{\Zitat}[2][\empty]{\ifthenelse{\equal{#1}{\empty}}{\parencite{#2}}{\parencite[#1]{#2}}}
\newcommand{\Autor}[1]{\textsc{#1}} % zum Ausgeben von Autoren
\newcommand{\itemd}[2]{\item{\textbf{#1}}\\{#2}} % erzeugt ein Listenelement mit fetter Überschrift

% fügt Tabellen aus einer TEX-Datei ein
\newcommand{\newTabelle}[3] % Parameter: caption, label, file
{\begin{table}[htb]
     \centering
     \singlespacing
     \input{table/#3}
     \caption{#1}
     \label{#2}
\end{table}}

\newcommand{\newTabelleAnhang}[1] % Parameter: file
{\begin{center}
     \singlespacing
     \input{table/#1}
\end{center}}

% einfaches Wechseln der Schrift, z.B.: \changefont{cmss}{sbc}{n}
\newcommand{\changefont}[3]{\fontfamily{#1} \fontseries{#2} \fontshape{#3} \selectfont}

% Verwendung analog zu \includegraphics
\newlength{\myx} % Variable zum Speichern der Bildbreite
\newlength{\myy} % Variable zum Speichern der Bildhöhe
\newcommand\includegraphicstotab[2][\relax]{%
% Abspeichern der Bildabmessungen
    \settowidth{\myx}{\includegraphics[{#1}]{#2}}%
    \settoheight{\myy}{\includegraphics[{#1}]{#2}}%
% das eigentliche Einfügen
    \parbox[c][1.1\myy][c]{\myx}{%
        \includegraphics[{#1}]{#2}}%
}


% Tabellen, die den Zwischenstand nach einer Projektphase verdeutlichen
\newcommand{\Zwischenstand}[2]{\subsection{Zwischenstand}Tabelle~\ref{tab:#2} zeigt den Zwischenstand nach der #1.\newTabelle{Zwischenstand nach der #1}{tab:#2}{#2.tex}}

% Verweise
\newcommand{\Verweis}[1]{Kapitel~\ref{#1}}
\newcommand{\VerweisKL}[1]{Kapitel~\ref{#1} auf Seite \pageref{#1}}
\newcommand{\VerweisVL}[1]{Kapitel~\ref{#1}: \nameref{#1} auf Seite \pageref{#1}}
\newcommand{\VerweisKT}[1]{(Siehe Kapitel~\ref{#1} auf Seite \pageref{#1})}
\newcommand{\VerweisVT}[1]{(Siehe Kapitel~\ref{#1}: \nameref{#1} auf Seite \pageref{#1})}
\newcommand{\VerweisKF}[1]{\footnote{Siehe Kapitel~\ref{#1} auf Seite \pageref{#1}.}}
\newcommand{\VerweisVF}[1]{\footnote{Siehe Kapitel~\ref{#1}: \nameref{#1} auf Seite \pageref{#1}.}}

% ### Semantische Textverarbeitung ###

% Abkürzung
\newcommand{\Abkuerzung}[1]{\textit{\ac{#1}}}

% Begriffe
\newcommand{\Fachbegriff}[1]{\textbf{#1}}
\newcommand{\Klasse}[1]{\texttt{#1}}
\newcommand{\Methode}[1]{\texttt{#1}}
\newcommand{\MethodenRef}[1]{\texttt{#1}}
\newcommand{\Code}[1]{\textsf{#1}}



% ----- Start document --------------------------
\begin{document}

% ----- cover -----------------------------------
    \clearpage

    \thispagestyle{plain}
\begin{titlepage}

    \onehalfspacing

    \begin{center}
        \includegraphics[scale=0.3]{picture/ihk_logo2.png}
        \footnote{\url{http://www.hfci.de/uploads/images/zoom/ihk-oldenburg-logo.png}}\\[1ex]
        \Large{Abschlussprüfung \pruefungstermin}\\[3ex]

        \Large{\ausbildungsberuf}\\
        \LARGE{\betreff}\\[4ex]
%\normalsize Abgabetermin: \ort, den \abgabetermin\\[4ex]

        \huge{\textbf{Erweiterung einer Webanwendung\\ um Task Notifications}}\\[4em]

        \normalsize
        {\setlength{\tabcolsep}{3em}
            \begin{tabular}{ll}
                \noindent\hspace*{13mm}\includegraphics[scale=0.55]{picture/scribble_mitarbeiter.png}
                &
                \includegraphics[scale=0.25]{picture/scribble_btc.png}
                \\
                \textbf{Prüfungsbewerber:}
                &
                \textbf{Ausbildungsbetrieb:}
                \\[1em]
                \autor
                &
                BTC AG
                \\
                Lamberti Str. 72
                &
                Escherweg 5
                \\
                26121 Oldenburg
                &
                26121 Oldenburg
                \\
                Geb.: 10.07.1993 in Bremen
                \\
                Azubi-Ident.-Nr.: 0001014762
                \\[4em]
            \end{tabular}
        }

    \end{center}

    \small
    \noindent
    Dieses Werk einschließlich seiner Teile ist \textbf{urheberrechtlich geschützt}.
    Jede Verwertung außerhalb der engen Grenzen des Urheberrechtgesetzes ist ohne
    Zustimmung des Autors unzulässig und strafbar. Das gilt insbesondere für
    Vervielfältigungen, Übersetzungen, Mikroverfilmungen sowie die Einspeicherung
    und Verarbeitung in elektronischen Systemen.

\end{titlepage}

% ----- preface ---------------------------------
    \clearpage
    \newpage\null\thispagestyle{empty}\newpage % empty page
    \pagenumbering{Roman}

    \begingroup
\hypersetup{linkcolor = black}

\pdfbookmark[1]{Inhaltsverzeichnis}{toc}
\addtocontents{toc}{\vspace{-14px}}
\tableofcontents
\addtocontents{toc}{\protect\setstretch{1.2}}
\clearpage

\listoftables
\clearpage

\listoffigures
\clearpage

\endgroup

\newcommand{\abkvz}{Abkürzungsverzeichnis}
\renewcommand{\nomname}{\abkvz}
\section*{\abkvz}
\label{toc:\abkvz}
\markboth{\abkvz}{\abkvz}
%\addcontentsline{toc}{section}{\abkvz}
% Es werden nur die Abkürzungen aufgelistet, die mit \ac definiert und auch benutzt wurden. 
%
% \acro{ABC}{Aaa Bbb Ccc\acroextra{ (Ddd)}}
% Ergibt in der Liste:
%  ABC Aaa Bbb Ccc (Ddd)
% Im Text aber:
%  \ac{ABC} -> Aaa Bbb Ccc (ABC)

\begin{acronym}
    \acro{API}{Application Programming Interface}
    \acro{IDE}{Integrated Development Environment}
    \acro{Java EE}{Java Enterprise Edition}
    \acro{JSON}{JavaScript Object Notation}
    \acro{SQL}{Structured Query Language}
    \acro{UI}{User Interface}
\end{acronym}



% ----- content ---------------------------------
    \clearpage
    \newpage\null\thispagestyle{empty}\newpage % empty page
    \pagenumbering{arabic}

    \section{Einleitung}
\label{sec:Einleitung}

\subsection{Vorwort}
\label{sub:Vorwort}
fghjfghjfghj TAI\footnote{Technischer Assistent für Informatik.} fghjfghj im Kapitel~\ref{sub:Projektumfeld} fghjfghj

\subsection{Kurzporträt der Ausbildungsstätte}
\label{sub:Kurzporträt der Ausbildungsstätte}
dfghdfghdfgh \Abkuerzung{UI} dfghdfgh

\begingroup
\setlength\intextsep{0pt}

\subsection{Projektumfeld}
\label{sub:Projektumfeld}
dfghdfghdfgh Non-SAP Technologien d dfghdfghdfgh Non-SAP Technologien d
dfghdfghdfgh Non-SAP Technologien d dfghdfghdfgh Non-SAP Technologien d
dfghdfghdfgh Non-SAP Technologien d dfghdfghdfgh Non-SAP Technologien d
dfghdfghdfgh Non-SAP Technologien d dfghdfghdfgh Non-SAP Technologien d

\begin{wrapfigure}[15]{r}[25pt]{0\textwidth}
    \includegraphicsKeepAspectRatio{./picture/systemlandschaft.pdf}{0.6}
    \caption{Systemlandschaft}
    \label{app:Systemlandschaft}
\end{wrapfigure}
sddfghdfgh \zB dfghdfgh (Siehe \Abbildung{app:Systemlandschaft})

dfghdfghdfgh Non-SAP Technologien d dfghdfghdfgh Non-SAP Technologien d
dfghdfghdfgh Non-SAP Technologien d dfghdfghdfgh Non-SAP Technologien d
dfghdfghdfgh Non-SAP Technologien d dfghdfghdfgh Non-SAP Technologien d
dfghdfghdfgh Non-SAP Technologien d dfghdfghdfgh Non-SAP Technologien d
dfghdfghdfgh Non-SAP Technologien d dfghdfghdfgh Non-SAP Technologien d
dfghdfghdfgh Non-SAP Technologien d dfghdfghdfgh Non-SAP Technologien d
dfghdfghdfgh Non-SAP Technologien d dfghdfghdfgh Non-SAP Technologien d
dfghdfghdfgh Non-SAP Technologien d dfghdfghdfgh Non-SAP Technologien d
dfghdfghdfgh Non-SAP Technologien d dfghdfghdfgh Non-SAP Technologien d
dfghdfghdfgh Non-SAP Technologien d dfghdfghdfgh Non-SAP Technologien d

\endgroup

\subsection{Motivation und Zielsetzung}
\label{sub:Motivation und Zielsetzung}
dfghdfgh

\subsection{Erklärung zum Dokument}
\label{sub:Erklärung zum Dokument}
An dieser Stelle wird die Gestaltung des Dokuments erläutert.
\begin{itemize}[noitemsep]
    \item Ein {[\textit{kursiver Text}]} weist auf eine Abkürzung hin, die im Abkürzungsverzeichnis zu finden ist.
    \item Ein {[\textbf{fettgedruckter Text}]}, mit Ausnahme von Überschriften, weist auf einen Fachbegriff hin, zu dem im Dokument
    eine Erklärung vorliegt. Ausgenommen sind Erklärungen in Fußnoten.
    \item Ein {[\texttt{Text in Schreibmaschinenschrift}]} weist auf eine Klasse oder Methode hin.
    \item Ein {[\textsf{Text in serifloser Schrift}]} weist auf sonstige Programmcode-Textbausteine hin.
    \item Liegt das Dokument in elektronischer Form vor, sind Abkürzungen und Quellenangaben zum jeweiligen Verzeichniseintrag
    verlinkt.
\end{itemize}

\subsection{Aufbau des Dokuments}
\label{sub:Aufbau des Dokuments}
dfghdfgh


    \section{Planungsphase}
\label{sec:Planungsphase}

\subsection{Zeitlicher Rahmen des Projektes}
\label{sub:Zeitlicher Rahmen des Projektes}
sdfgsdfgsdg

\subsection{Vorgehensweise}
\label{sub:Vorgehensweise}
sdfgsdfgsdfg

\subsection{Verwendete Hardware}
\label{sub:Verwendete Hardware}
sdfgsdfg \Tabelle{tab:Laptop-Spezifikation} auf S.~\pageref{app:Laptop-Spezifikation} sdfgsdfg 24''

\subsection{Verwendete Software/Technologien}
\label{sub:Verwendete Software/Technologien}
asdfasdfasdf
\begin{description}
    \item [Intellij IDEA] asdfasdfasdf (von engl.: \Abkuerzung{IDE}).
    \item [Java EE]
    Die Plattform \Abkuerzung{Java EE}, asdfasdf / \Abkuerzung{API}'s asdfasdf
    \item [Atmosphere]
    asdfasdfasdf Servlet\footnote{Easdfasdfasdf} \Abkuerzung{API} der asdfasdfasdf
    \item [Spring]
    asdfasdf Java/\Abkuerzung{Java EE} zasdfasdfa
    \item [jQuery]
    asdfasdf
\end{description}
    \section{Analysephase}
\label{sec:Analysephase}
dfghdfgh

\subsection{Soll-Konzept}
\label{sub:Soll-Konzept}
sdfsdfsdf

\subsection{Ist-Analyse}
\label{sub:Ist-Analyse}
dfghdfg
asdfasdf \AbbildungV{app:Mockup} asdfgasdf

\begin{figure}[t]
    \centering
    \includegraphicsKeepAspectRatio{picture/mockup_status-anzeige.png}{1}
    \caption{Mockup Status-Anzeige - Einbettung und Verknüpfung zur Task-Übersicht}
    \label{app:Mockup}
\end{figure}

\subsubsection{Kosten-Analyse}
\label{subsub:Kosten-Analyse}
sdfgsdfgsdfg:
\[ 70h * 45{\frac{\eur{}}{h}} = \eur{3150} \]

\subsubsection{Nutzen-Analyse}
\label{subsub:Nutzen-Analyse}
asdfasdf "`\titel"' asdfasdf

\paragraph{Monetärer Vorteil}
asdfasdf \eur{145} asdfasdfasdf:
\[ \bigg(70h * 145{\frac{\eur{}}{h}}\bigg) - \eur{3150} = \eur{7000} \]


    \section{Konzeptphase}
\label{sec:Konzeptphase}
sdfgsdfg . sdfgsdfg \Abkuerzung{UI}-Entwurf sdfgsdfg

\subsection{Technisches Konzept}
\label{sub:Technisches Konzept}
sdfgsdfg

\begin{description}
    \item [Kommunikation: Server zu Client]
    Das Framework \Fachbegriff{Atmosphere} sdfgsdfg \zB sdfgsdfgsdfg \sog "`Heartbeats"' (zu Deutsch: Herzschläge) sdfgsdfgsd
    \item [Datenhaltung]
    sdfgsdfg \VerweisKT{subsub:UserStatusRepository} sdfgsdfg
    \item [Datenübertragung]
    sdfgsdfg \Abkuerzung{JSON}\footnote{Erklärung siehe .} sdfgsdfgsdfg
\end{description}

\subsection{UI-Entwurf}
\label{sub:UI-Entwurf}
sdfgsdfgsdfg\\
sdfgsdfgsdfg \Anhang{app:UI-Entwurf} sdfgsdfgdsgf \VerweisKL{sub:Soll-Konzept} sdfgsdfgsdfg

\subsection{Proof-of-Concept}
\label{sub:Proof-of-Concept}
sdfgsdfgsdfg

\subsection{Testplan}
\label{sub:Testplan}
sdfgsdfgsdfg \TabelleK{tab:Testfälle} sdfgsdfgsdfg

\newTabelle{Testfälle}{tab:Testfälle}{Testfälle.tex}
    \clearpage


\section{Realisierungsphase}
\label{sec:Realisierungsphase}
asdfasdf

\subsection{Projektstruktur}
\label{sub:Projektstruktur}
\footnotetext{Internet Explorer (bis Version 7 pflicht), Mozilla Firefox, Google Chrome, Opera.}
\begin{wrapfigure}[14]{r}[25pt]{0\textwidth}
    \tikzstyle{every node}=[draw=black,thick,anchor=west]
    \tikzstyle{new}=[dashed,fill=hig]
    \scalebox{0.8}{
        \begin{tikzpicture}
            [%
            grow via three points={one child at (0.5,-0.7) and
            two children at (0.5,-0.7) and (0.5,-1.4)},
            edge from parent path={(\tikzparentnode.south) |- (\tikzchildnode.west)}]
            \texttt{
                \node {com.<mandant>.contrax}
                child { node {party.domain}
                child { node [new] {repository}
                }
                }
                child [missing] {}
                child { node {shared.infrastructure.async}
                child { node [new] {atmosphere}
                child { node [new] {domain}
                child { node [new] {model}
                }
                child { node [new] {service}
                child { node [new] {impl}
                }
                }
                }
                }
                };
            }
        \end{tikzpicture}
    }
    \caption{Differenz zur vorherigen Packetstruktur}
    \label{app:DifferenzPacketstruktur}
\end{wrapfigure}
asdfasdf "`Domain Driven
Design"'\footnote{Definition siehe .} asdfasdf \Code{"`com.<mandant>.contrax"'} asdfasdf \Code{<mandant>}
asdfasdfasdf \AnhangK{app:Paketstruktur} asdfasdf \Abbildung{app:DifferenzPacketstruktur} asdfasdfasdf

\subsection{Einrichten/Konfigurieren von Atmosphere}
\label{sub:Einrichten/Konfigurieren von Atmosphere}
Dasdfasdf \Anhang{app:pom.xml} asdfasdf Maven\footnote{Maven ist ein "`Build-Management-Tool"' basierend auf
Java.} asdfasdf \Klasse{VmsWebAppInitializer} asdfasdf \Code{web.xml} asdfasdf \Code{"/notify/*"} asdfasdf
\Fachbegriff{Spring} asdfasdf (siehe )

\subsection{Kommunikation: Server zu Client}
\label{sub:Kommunikation: Server zu Client}
asdfasdf asdf \Anhang{app:main.js}
\begin{description}
    \item [AtmosphereHandler] asdfasdf \Klasse{AtmosphereHandler} \Code{"/notify/*"} asdfas
    \item [Broadcaster] asdfasdf \Methode{broadcast} asdfasdf \MethodenRef{AtmosphereResource\#getBroadcaster()} asdf
    aasdfasfd \sog \Code{Factory}\footnote{asdfasdf \Code{Factory}-asdfasdfe \zB asdfasdf} asdf
\end{description}

\subsubsection{Verbindungshaltung}
\label{subsub:Verbindungshaltung}
asdfasdf (\Code{\acs{SQL}HttpRequest}) asdfasdf \Fachbegriff{Resource} asdfasdf \VerweisVL{sub:Soll-Konzept} aufgezeigt.\\
asdfasdf

\subsection{Datenhaltung}
\label{sub:Datenhaltung}
asdfasdfasdf \AnhangKapitel{app:Quellcode}

\subsubsection{UserStatusRepository}
\label{subsub:UserStatusRepository}
\begin{description}
    \item [Repository]
    asdfasdf
\end{description}

Zum Zeitpunkt der Implementierung wird bisher nur der letzte Besuch der Offenen-Task-Übersicht eines jeden Benutzers im
\Fachbegriff{Repository} gespeichert. Die Ermittlung des Zeitstempels wird unten stehend erklärt. Da es sich lediglich um ein
\Code{DateTime}-Objekt handelt, könnte diese Information in der vorhandenen Datenbank gespeichert werden. Gemäß dem Abschnitt
Datenhaltung aus \VerweisVL{sub:Technisches Konzept} wurde jedoch das \Fachbegriff{Repository} entwickelt.

\subsubsection{TaskNotificationStatisticsModel}
\label{subsub:TaskNotificationStatisticsModel}
asdfasdf \VerweisVF{sub:Soll-Konzept} asdfasdf

\subsection{Datenübertragung}
\label{sub:Datenübertragung}
asdfasdf \ggfs asdfasdf

\subsection{UI-Entwurfsumsetzung}
\label{sub:UI-Entwurfsumsetzung}
asdfasdf \VerweisKL{sub:Soll-Konzept} sdfg


    \section{Qualitätssicherungsphase}
\label{sec:Qualitätssicherungsphase}
asdfasdf

\subsection{Definition of Done}
\label{sub:Definition of Done}
aasdfasdf\footnote{Vgl. \url{http://www.clean-code-developer.de/Roter-Grad.ashx\#Praktiken_10} (Aufruf: 02.04.2015).} asdfasdf

"`asdfasdf"'

asdfasdf

\begingroup
\setlength\intextsep{0pt}
\begin{wraptable}[4]{r}[-24pt]{0\textwidth}

    \begin{tabular}{rl|c}

        \hline
        \rowcolor{head} \multicolumn{1}{c}{\textbf{Nr.}} & \multicolumn{1}{c|}{\textbf{Testfall}} &
        \multicolumn{1}{|c}{\textbf{Auswertung}} \\
        \hline
        1               & Consistent\linebreak \Fachbegriff{Repository} & \checkmark \\
        \hline
        \rowcolor{odd}2 & Task statistics validated                     & \checkmark \\
        \hline
        3               & Connectivity                                  & \checkmark \\
        \hline
        \rowcolor{odd}4 & Incoming\linebreak event                      & \checkmark \\
        \hline
        5               & Data transfer                                 & \checkmark \\
        \hline
        \rowcolor{odd}6 & Frontend event handling                       & \checkmark \\
        \hline
        7               & User scoping                                  & \checkmark \\
        \hline
        \rowcolor{odd}8 & Accessibility                                 & \checkmark \\
        \hline
        9               & Browser\linebreak compatibility               & \checkmark \\

    \end{tabular}

    \caption{Testergebnisse}
    \label{tab:Testergebnisse}

\end{wraptable}


% \begin{tabular}{rp{2.5cm}|p{5.5cm}|c}
% 
% \hline
% \multicolumn{1}{c}{\textbf{Nr.}}   & \multicolumn{1}{c|}{\textbf{Testfall}} & \multicolumn{1}{c}{\textbf{Erwartetes Ergebnis}} &
% \multicolumn{1}{|c}{\textbf{Auswertung}} \\
% \hline
% 1     & Consistent\linebreak \Fachbegriff{Repository} & Daten sind vorhanden und können nur vom \Fachbegriff{Repository}
% modifiziert werden & \checkmark\\
% \hline
% \rowcolor{odd}2     & Task statistics validated & Korrektes Ergebnis der Tasks-Statistik wird geliefert & \checkmark\\
% \hline
% 3     & Connectivity & Client-Server Verbindung besteht (\acs{XML}HttpRequest auf "pending") & \checkmark\\
% \hline
% \rowcolor{odd}4     & Event\linebreak incoming & Im Frontend wird die JavaScript-Methode \Methode{onMessage} angesprochenn &
% \checkmark\\
% \hline
% 5     & Data transfer & Event-Daten sind im \Abkuerzung{JSON}-Format & \checkmark\\
% \hline
% \rowcolor{odd}6     & Frontend event handling & Task-Event sorgt für Aktualisierung der \Abkuerzung{UI} & \checkmark\\
% \hline
% 7     & User scoping & Nur gewünschte Benutzer werden benachrichtigt & \checkmark\\
% \hline
% \rowcolor{odd}8     & Accessibility  & Status-Anzeigen sind barrierefrei Nutzbar & \checkmark\\
% \hline
% 9     & Browser\linebreak compatibility & Das Design funktioniert in allen Browser angemessen  & \checkmark\\
% 
% \end{tabular}

\subsection{Testergebnisse}
\label{sub:Testergebnisse}
asdfasdf\Tabelle{tab:Testergebnisse} asdfasdf \TabelleV{tab:Testfälle} asdfasdf

    \section{Abschlussphase}
\label{sec:Abschlussphase}
sdfgsdfgsdfgsdfg

\subsection{Soll-Ist-Vergleich}
\label{sub:Soll-Ist-Vergleich}
sdfgsdfgsdfg \\
\endgroup
In \TabelleK{tab:Zeitplanung} sdfgsdfg

\newTabelle{sdfgsdfgsdfgsdfgsdfgsdfgsdfgsdfgsdfg}{tab:Zeitplanung}{Zeitplanung.tex}
sdfgsdfgsdfg


\parencite{CCCCC} asdfasdf
\cite{CCCCC} asdfasdfasdf

\parencite{BBBBB} asdfasdf
\cite{BBBBB} asdfasdfasdf
\parencite[vgl.][aa]{BBBBB} asdfasdfasdf
\parencite[vgl.]{BBBBB} asdfasdfasdf

\parencite{AAAAA} asdfasdf
\cite{AAAAA} asdfasdfasdf


% ----- declaration on oath ---------------------
    \clearpage
    \pagenumbering{arabic}

    \addsec{Eidesstattliche Erklärung}

Ich, \autor, versichere hiermit, dass ich meine \textbf{\betreff} mit dem Thema
\begin{quote}
    \textit{\kompletterTitel}
\end{quote}
selbständig verfasst und keine anderen als die angegebenen Quellen und Hilfsmittel benutzt habe,
wobei ich alle wörtlichen und sinngemäßen Zitate als solche gekennzeichnet habe. Die Arbeit
wurde bisher keiner anderen Prüfungsbehörde vorgelegt und auch nicht veröffentlicht.\\[6ex]

\ort, den \abgabetermin


\rule[-0.2cm]{5.5cm}{0.5pt}

\textsc{\autor} 


% ----- bibliography ----------------------------
    \clearpage
    \pagenumbering{roman}

    \renewcommand{\refname}{Literaturverzeichnis}
    \printbibliography
% TODO may divide sources, see https://www.overleaf.com/learn/latex/Articles/Getting_started_with_BibLaTeX

% ----- appendix --------------------------------
    \clearpage
    \pagenumbering{roman}

    \appendix
\singlespacing
\addtocontents{toc}{\protect\setcounter{tocdepth}{3}}


\section{Anhang}

\subsection{Weboberfläche}
\label{app:Weboberfläche}

\subsubsection{UI-Entwurf}
\label{app:UI-Entwurf}
asdfasdfasdf
\begin{figure}[htb]
    \centering
    \includegraphicsKeepAspectRatio{appendix/ui_entwurf.png}{1.0}
    \caption{UI-Entwurf}
\end{figure}

\subsubsection{listOpenTasks.xhtml (Auszug)}
\label{app:listOpenTasks}
\lstinputlisting[language=HTML]{appendix/listOpenTasks.xhtml}

\clearpage

\subsubsection{main.less (Auszug)}
\label{app:main.less}
\lstinputlisting[language=CSS]{appendix/main.css}


\clearpage

\subsection{Tabellen}
\label{Tabellen}

\subsubsection{Laptop-Spezifikation}
\label{app:Laptop-Spezifikation}
\newTabelle{Laptop-Spezifikation}{tab:Laptop-Spezifikation}{Laptop.tex}

\subsubsection{Auszug der vorhandene Paketstruktur}
\label{app:Paketstruktur}
asdfasdf: \Code{com.contract.domain.model.transactions}

\clearpage

\subsection{Quellcode}
\label{app:Quellcode}
asdfasdf

\subsubsection{AbstractNotificationService.java}
\label{app:AbstractNotificationService}
\lstinputlisting[style=Java]{appendix/AbstractNotificationService.java}

\subsubsection{main.js (Auszug)}
\label{app:main.js}
asdfasdfa \\
\Methode{configureListeners(request,socket)} asdfasdfasdf
\lstinputlisting[language=JavaScript]{appendix/main.js}

\subsubsection{pom.xml (Auszug)}
\label{app:pom.xml}
\lstinputlisting[style=XML]{appendix/pompom.xml}

\clearpage

\subsection{Weitere Grafiken}
\label{app:Grafiken}
dfghdfgh



% ----- End document ----------------------------
\end{document}

